\documentclass[conference]{IEEEtran}

\usepackage{graphicx}
\usepackage{amsmath,amssymb}
\usepackage{float}
\usepackage{booktabs}
\usepackage{multirow}

\begin{document}

\title{
    \textbf{Benchmark Report: 508.namd\_r (Train Input)} \\
    \text{Register Spill Analysis on RISC-V via gem5}
}

\author{
    \IEEEauthorblockN{Lütfullah Burak Kaya}
    \IEEEauthorblockA{
        Ozyegin University \\
        burak.kaya.50711@ozu.edu.tr
    }
    \and
    \IEEEauthorblockN{Ismail Akturk}
    \IEEEauthorblockA{
        Ozyegin University \\
        ismail.akturk@ozyegin.edu.tr
    }
}

\newcommand{\LongDate}{%
    \number\day\space
    \ifcase\month
    \or January\or February\or March\or April\or May\or June%
    \or July\or August\or September\or October\or November\or December%
    \fi
    \space \number\year
}

\twocolumn[
\begin{@twocolumnfalse}
\maketitle
\begin{center}{\small \LongDate}\end{center}
\vspace{1em}
\end{@twocolumnfalse}
]

% ================================================

\begin{abstract}
This report presents the register spill characterization of the
\texttt{508.namd\_r} benchmark from SPEC CPU2017, executed on a
RISC-V (RV64GC) target using the gem5 TimingSimpleCPU simulator
with the \textbf{train} input dataset. A custom SpillDetector was
integrated into gem5 to count store-then-load patterns on the stack
within a user-defined Region of Interest (ROI). The train dataset
uses the same \texttt{apoa1.input} molecular dynamics system as the
test dataset but runs for \textbf{7 iterations} instead of 1,
yielding 218.4 billion ROI instructions. The measured spill rate is
\textbf{1.0464\%}, nearly identical to the test input rate of
1.0269\%, confirming that \texttt{namd\_r} exhibits highly stable,
input-size-independent register pressure. This places \texttt{namd\_r}
last among all 17 evaluated SPEC CPU2017 benchmarks in spill rate,
reflecting the benchmark's compute-intensive, numerically regular
character.
\end{abstract}

% =========================================================
\section{Benchmark Overview}

\texttt{508.namd\_r} is the SPEC CPU2017 rate variant of NAMD, a
parallel molecular dynamics application designed for high-performance
simulation of large biomolecular systems. The SPEC benchmark variant
runs serially on a single thread using a fixed input system
(\texttt{apoa1.input}), simulating the ApoA1 lipoprotein particle.

NAMD's computational kernel consists of tight numerical loops
evaluating pairwise forces between atoms. Each iteration of the main
loop executes a fixed sequence of floating-point force evaluations
(\texttt{RUN\_AND\_CHECKSUM}), result verification
(\texttt{WRITE}/\texttt{compare}), and atom position updates
(\texttt{moveatoms}). The regularity of this loop structure limits
the live variable count at any point: at most a handful of accumulators
and array pointers are simultaneously live, keeping register pressure
uniformly low across all iterations.

The RISC-V binary (\texttt{namd\_r.riscv}) was compiled with gem5
ROI markers inserted around the main iteration for-loop in
\texttt{src/spec\_namd.C}:
\begin{itemize}
    \item \texttt{m5\_work\_begin(0,0)} --- before \texttt{for (i = 0; i < iterations; ++i)}
    \item \texttt{m5\_work\_end(0,0)}   --- after the closing brace of the loop
\end{itemize}
This ROI encompasses all force-evaluation, checksum, and atom-movement
computation.

\subsection{Input Datasets}

SPEC CPU2017 provides three dataset sizes for \texttt{namd\_r}
(Table~\ref{tab:inputs}). All three use the same \texttt{apoa1.input}
file; the only difference is the iteration count.

\begin{table}[H]
\centering
\caption{namd\_r Input Dataset Sizes}
\label{tab:inputs}
\begin{tabular}{lll}
\toprule
\textbf{Dataset} & \textbf{Input file} & \textbf{Iterations} \\
\midrule
test  & \texttt{apoa1.input} & 1 \\
train & \texttt{apoa1.input} & 7 \\
ref   & \texttt{apoa1.input} & 38 \\
\bottomrule
\end{tabular}
\end{table}

The \textbf{train} dataset (7 iterations) was used for this study.
A single gem5 invocation was run with the \texttt{namd.in} argument
file specifying \texttt{--iterations 7}.

% =========================================================
\section{Simulation Setup}

\subsection{gem5 Configuration}

\begin{itemize}
    \item \textbf{CPU model:}   TimingSimpleCPU
    \item \textbf{ISA:}         RISC-V RV64GC
    \item \textbf{Caches:}      enabled (default L1 I/D + L2)
    \item \textbf{Memory:}      4\,GB simulated physical RAM
    \item \textbf{Spill mode:}  summary only (\texttt{spill\_verbose=False})
\end{itemize}

A single gem5 invocation was launched:

\begin{verbatim}
./build/RISCV/gem5.opt
  --outdir=benchmarks/namd/m5out_train
  configs/deprecated/example/se.py
  --cmd=benchmarks/namd/namd_r.riscv
  --options="--input benchmarks/namd/apoa1.input
    --iterations 7
    --output benchmarks/namd/apoa1.train.output"
  --cpu-type=TimingSimpleCPU
  --caches --mem-size=4GB
\end{verbatim}

% =========================================================
\section{Simulation Results}

\subsection{Pre-ROI Context}

Table~\ref{tab:preroi} shows the global counters at ROI entry,
reflecting the benchmark's initialisation and input-parsing phase.

\begin{table}[H]
\centering
\caption{Global Counters at ROI Entry}
\label{tab:preroi}
\begin{tabular}{lr}
\toprule
\textbf{Metric} & \textbf{Value} \\
\midrule
Total instructions (pre-ROI) & 1,558,178,337 \\
Total spills detected (pre-ROI) & 125,215,611 \\
Pre-ROI spill rate & 8.04\% \\
\bottomrule
\end{tabular}
\end{table}

The notably higher pre-ROI spill rate (8.04\%) relative to the ROI
spill rate (1.05\%) reflects the initialisation code paths: NAMD's
startup reads and parses the input structure file, constructs
patch/atom data structures, and sets up pairlists --- operations that
invoke many function calls with rich call chains and consequently
higher register pressure than the steady-state force-evaluation loop.

\subsection{ROI Spill Statistics}

Table~\ref{tab:roi} presents the complete ROI statistics for the
train run.

\begin{table}[H]
\centering
\caption{ROI Spill Statistics --- 508.namd\_r (train)}
\label{tab:roi}
\begin{tabular}{lr}
\toprule
\textbf{Metric} & \textbf{Value} \\
\midrule
ROI instructions   & 218,427,292,249 \\
ROI stores         &  13,245,362,832 \\
ROI loads          &  57,662,937,316 \\
\midrule
\textbf{ROI spills} & \textbf{2,285,717,035} \\
\textbf{Spill rate} & \textbf{1.0464\%} \\
\bottomrule
\end{tabular}
\end{table}

\subsection{Timing}

\begin{table}[H]
\centering
\caption{Simulation Timing}
\label{tab:timing}
\begin{tabular}{lrr}
\toprule
\textbf{Dataset} & \textbf{Sim time} & \textbf{Host time} \\
\midrule
test (1 iter) &  68.73\,s &  25,169\,s ($\approx$7.0\,h) \\
train (7 iter) & 452.37\,s & 170,069\,s ($\approx$47.2\,h) \\
\bottomrule
\end{tabular}
\end{table}

Simulated time scales nearly linearly with iteration count
($452.37 / 68.73 \approx 6.6\times$ for $7\times$ the iterations),
confirming that each iteration performs an identical amount of work
with no warmup or convergence effect.

% =========================================================
\section{Analysis}

\subsection{Spill Rate Stability Across Input Sizes}

The train spill rate of \textbf{1.0464\%} is within 1.9\% relative
of the test rate (1.0269\%), despite a $6.8\times$ increase in ROI
instructions. This near-perfect stability confirms that NAMD's
register pressure is determined entirely by the structure of its
force-evaluation inner loop, which is identical across iterations.
Unlike compiler or interpreter workloads whose register pressure
varies with input content, NAMD executes the same instruction
sequence repeatedly with only the data values changing.

\subsection{Train vs.\ Test Input Comparison}

Table~\ref{tab:compare_test_train} summarises the test and train
results side by side.

\begin{table}[H]
\centering
\caption{namd\_r Test vs.\ Train Comparison}
\label{tab:compare_test_train}
\begin{tabular}{lrr}
\toprule
\textbf{Metric} & \textbf{Test (1 iter)} & \textbf{Train (7 iter)} \\
\midrule
ROI instructions & 31.9\,B  & 218.4\,B \\
ROI stores       &  1.9\,B  &  13.2\,B \\
ROI loads        &  8.4\,B  &  57.7\,B \\
ROI spills       & 327.7\,M &   2,285.7\,M \\
Spill rate       & 1.0269\% & \textbf{1.0464\%} \\
\bottomrule
\end{tabular}
\end{table}

All four counters scale by approximately $6.8\times$ (reflecting
7 iterations with minor overhead), while the spill rate remains
essentially constant. This validates the test run as a faithful
representative of NAMD's register-spilling behaviour.

\subsection{Load-to-Store Ratio}

The combined ROI load count ($57.7\,\text{B}$) is $4.35\times$
the store count ($13.2\,\text{B}$):
\[
\frac{57{,}662{,}937{,}316}{13{,}245{,}362{,}832} \approx 4.35
\]
This high ratio reflects the force-evaluation kernel's access pattern:
each atom pair evaluation reads positions, charges, and force
parameters from multiple arrays (many loads per atom pair) and writes
back only the accumulated force vector (few stores). The $4.35\times$
load-to-store ratio is the highest in the benchmark suite, consistent
with NAMD's streaming, read-dominated memory access pattern.

\subsection{Spill Fraction of Loads}

Within the ROI, the ratio of spills to total loads is:
\[
\frac{2{,}285{,}717{,}035}{57{,}662{,}937{,}316} \approx 3.96\%
\]
Approximately 1 in 25 load instructions is a spill reload. This is
the lowest spill-to-load fraction in the study, corroborating the
low absolute spill rate: NAMD's tight loop fits comfortably within
RISC-V's 32-register file for the majority of execution.

\subsection{ROI Coverage}

The ROI covers $218{,}427{,}292{,}249$ out of
$219{,}985{,}470{,}586$ total simulated instructions:
\[
\frac{218{,}427{,}292{,}249}{219{,}985{,}470{,}586} \approx 99.29\%
\]
The 0.71\% pre-ROI share corresponds to the 1.56\,B startup
instructions, which remain constant regardless of iteration count.
As iteration count increases further (e.g., the ref dataset at 38
iterations), ROI coverage approaches 100\%.

\subsection{Cross-Benchmark Context}

Table~\ref{tab:compare} places \texttt{namd\_r} (train) in the
context of all 17 evaluated SPEC CPU2017 benchmarks.

\begin{table}[H]
\centering
\caption{Cross-Benchmark Spill Rate Comparison}
\label{tab:compare}
\begin{tabular}{lrr}
\toprule
\textbf{Benchmark} & \textbf{ROI Spills} & \textbf{Spill Rate} \\
\midrule
502.gcc\_r (train)  & 17,918,365,207 & 8.1160\% \\
507.cactuBSSN\_r    &  3,782,099,438 & 7.8756\% \\
511.povray\_r       &    174,359,922 & 7.7809\% \\
500.perlbench\_r    &      1,070,055 & 7.2838\% \\
520.omnetpp\_r      &    894,619,681 & 7.1994\% \\
502.gcc\_r (test)   &      1,108,838 & 6.9185\% \\
526.blender\_r      &     61,845,869 & 6.2193\% \\
523.xalancbmk\_r    &     18,381,250 & 4.6427\% \\
541.leela\_r        &  1,158,236,271 & 4.4438\% \\
519.lbm\_r          &    274,130,440 & 4.2130\% \\
531.deepsjeng\_r    &  1,928,903,656 & 4.1749\% \\
525.x264\_r         &    767,981,981 & 3.2812\% \\
544.nab\_r          &    151,310,616 & 2.2612\% \\
538.imagick\_r      &      2,424,899 & 2.0601\% \\
557.xz\_r           &     38,411,428 & 1.9069\% \\
505.mcf\_r          &    446,113,136 & 1.8263\% \\
\textbf{508.namd\_r (train)} & \textbf{2,285,717,035} & \textbf{1.0464\%} \\
508.namd\_r (test)  &    327,718,955 & 1.0269\% \\
510.parest\_r       &    127,688,196 & 0.3427\% \\
\bottomrule
\end{tabular}
\end{table}

\texttt{namd\_r} ranks last among all 17 benchmarks, with a spill
rate 7.8$\times$ lower than \texttt{gcc\_r} (train) and only slightly
above \texttt{parest\_r} (0.34\%). The benchmark's scientific
computing character --- tight, numerically intensive loops with
regular memory access --- yields the minimum register pressure
observable in the SPEC CPU2017 C/C++ suite on RISC-V RV64GC.

% =========================================================
\section{Conclusion}

The \texttt{508.namd\_r} benchmark with the train input (7 iterations)
exhibits a register spill rate of \textbf{1.0464\%} across
218.4 billion ROI instructions, the lowest of all 17 evaluated SPEC
CPU2017 benchmarks. The rate is nearly identical to the test input
result (1.0269\%), confirming that NAMD's register pressure is
structurally determined by its force-evaluation loop and independent
of iteration count. The $4.35\times$ load-to-store ratio reflects the
benchmark's streaming read-dominated access pattern, and the 3.96\%
spill-to-load fraction is the lowest in the study. With 99.29\% ROI
coverage, the train run provides a thorough characterisation of
NAMD's register behaviour at scale. These results establish
\texttt{namd\_r} as the least register-pressure-intensive workload in
the SPEC CPU2017 suite on RISC-V RV64GC, making it a useful baseline
for studies requiring a low-spill-rate numerical benchmark.

\end{document}
